\documentclass[11pt,a4paper]{article}

% ---------- Preamble ----------
\usepackage[margin=25mm]{geometry}
\usepackage{amsmath,amssymb}
\usepackage{bm}
\usepackage{physics}
\usepackage{siunitx}
\usepackage{enumitem}

\title{Electromagnetic Force and Torque in a DC Motor}
\author{}
\date{}

\begin{document}

\maketitle

\section{Force on a Current-Carrying Conductor}

A conductor carrying an electric current in a magnetic field experiences
an electromagnetic force.

In vector form,

\[
\boxed{
\bm{F}
=
I\,\bm{\ell}\times\bm{B}
}
\]

where

\begin{itemize}
    \item $\bm{F}$ is the electromagnetic force,
    \item $I$ is the electric current,
    \item $\bm{\ell}$ is the length vector of the conductor,
    \item $\bm{B}$ is the magnetic flux density.
\end{itemize}

The current $I$ is a \emph{scalar}.  
The direction of the current is represented by the direction of
$\bm{\ell}$.

The magnitude of the force is therefore

\[
F = I\ell B\sin\theta,
\]

where $\theta$ is the angle between $\bm{\ell}$ and $\bm{B}$.

If the conductor is perpendicular to the magnetic field,

\[
\theta = 90^\circ,
\]

and hence

\[
\boxed{
F = BI\ell
}
\]

since $\sin 90^\circ = 1$.

\section{From Force to Torque}

Torque is produced when a force acts at a distance from the axis of
rotation.

In vector form,

\[
\boxed{
\bm{T}
=
\bm{r}\times\bm{F}
}
\]

where $\bm{r}$ is the position vector from the axis of rotation to the
point at which the force acts.

Its magnitude is

\[
T = rF\sin\alpha.
\]

For the usual geometry of a DC motor, the electromagnetic force acts
approximately perpendicular to the radius, so

\[
T = rF.
\]

Using

\[
F = BI\ell,
\]

we obtain

\[
T = rBI\ell.
\]

Thus, for a motor of fixed geometry,

\[
\boxed{
T \propto BI
}
\]

because $r$ and $\ell$, together with the number and arrangement of
conductors, are fixed by the construction of the motor.

\section{Magnetic Flux and Magnetic Flux Density}

For a uniform magnetic field perpendicular to an area $A$,

\[
\boxed{
\phi = BA
}
\]

where

\begin{itemize}
    \item $\phi$ is the magnetic flux,
    \item $B$ is the magnetic flux density,
    \item $A$ is the area through which the magnetic field passes.
\end{itemize}

For a motor of fixed geometry, $A$ is constant. Therefore,

\[
B \propto \phi.
\]

Substituting this relationship into

\[
T \propto BI,
\]

gives

\[
\boxed{
T \propto \phi I
}
\]

and for the armature of a DC motor,

\[
\boxed{
T = K_T \phi I_a
}
\]

where

\begin{itemize}
    \item $T$ is the electromagnetic torque,
    \item $K_T$ is the torque constant,
    \item $\phi$ is the magnetic flux per pole,
    \item $I_a$ is the armature current.
\end{itemize}

Therefore, the electromagnetic torque of a DC motor is proportional to
both the magnetic flux and the armature current.

\section{Back EMF}

When the armature rotates in a magnetic field, an electromotive force
is induced in the armature conductors.

This induced voltage opposes the applied voltage and is therefore
called the \emph{back electromotive force} (back EMF).

For a DC motor,

\[
\boxed{
E = K_E \phi N
}
\]

where

\begin{itemize}
    \item $E$ is the back EMF,
    \item $K_E$ is the EMF constant,
    \item $\phi$ is the magnetic flux per pole,
    \item $N$ is the rotational speed.
\end{itemize}

Hence,

\[
\boxed{
E \propto \phi N
}
\]

For constant magnetic flux,

\[
\boxed{
E \propto N
}
\]

so the back EMF increases as the motor rotates faster.

\section{Armature Voltage Equation}

Applying Kirchhoff's voltage law to the armature circuit gives

\[
\boxed{
V = E + I_a R_a
}
\]

or equivalently,

\[
\boxed{
E = V - I_a R_a
}
\]

where

\begin{itemize}
    \item $V$ is the terminal voltage,
    \item $E$ is the back EMF,
    \item $I_a$ is the armature current,
    \item $R_a$ is the armature resistance.
\end{itemize}

Thus,

\[
\boxed{
I_a = \frac{V-E}{R_a}
}
\]

The term $I_aR_a$ is an actual resistive voltage drop, whereas $E$ is
an induced voltage caused by the rotation of the armature.

\section{Mechanical Power}

For translational motion,

\[
W = Fx.
\]

For rotational motion, the corresponding relation is

\[
\boxed{
W = T\theta
}
\]

where $\theta$ is the angular displacement.

Dividing by time,

\[
P
=
\frac{W}{t}
=
T\frac{\theta}{t}.
\]

Since angular velocity is

\[
\omega = \frac{\theta}{t},
\]

the mechanical power is

\[
\boxed{
P = T\omega
}
\]

Thus, the correspondence between translational and rotational motion is

\[
\boxed{
\begin{aligned}
W &= Fx
&\quad&\longleftrightarrow&
W &= T\theta,\\
P &= Fv
&&\longleftrightarrow&
P &= T\omega.
\end{aligned}
}
\]

\section{Summary}

The important relationships for a DC motor are

\[
\boxed{
\bm{F}=I\bm{\ell}\times\bm{B}
}
\]

\[
\boxed{
T=K_T\phi I_a
}
\]

\[
\boxed{
E=K_E\phi N
}
\]

\[
\boxed{
V=E+I_aR_a
}
\]

and

\[
\boxed{
P=T\omega.
}
\]

The physical chain can be summarized as

\[
\text{Magnetic field + armature current}
\longrightarrow
\text{electromagnetic force}
\longrightarrow
\text{torque}
\longrightarrow
\text{rotation}.
\]

At the same time, rotation of the armature in the magnetic field
produces the back EMF:

\[
\text{magnetic flux + rotation}
\longrightarrow
\text{back EMF}.
\]

\end{document}
